\section{Spr6 Instance Inventory}
\label{app:spr6}

The Spr6 session is the instance behind every headless experiment in Section~\ref{sec:experiments}
and Section~\ref{sec:rl}. It is recorded separately from the current dataset because its counts
differ, and because quoting a conflict rate without the instance it came from was the single
largest source of confusion during the phase.

\subsection{The solver instance}

\begin{table}[H]
\centering
\caption{Properties of \texttt{rl/instances/spr2026.xml}.}
\label{tab:spr6inst}
\footnotesize
\begin{tabular}{@{}L{5.6cm}L{8.6cm}@{}}
\toprule
\textbf{Property} & \textbf{Value} \\
\midrule
Format            & CPSolver solver-model XML, \texttt{<timetable version="2.5">} \\
Origin            & \emph{Export XML} on the UniTime course timetabling solver screen \\
Session           & Spr6, 2026 \\
Classes (variables) & 262 \\
Students          & 990, carrying 7{,}748 enrolments \\
Rooms indexed     & 44 \\
Initial conflict rate & $C = 0.5313$ (526 of 990 students) \\
Self-contained    & Yes: rooms, instructors, time and room domains, distribution constraints,
                    students and the current assignment are all inside the file \\
\bottomrule
\end{tabular}
\end{table}

This format is \emph{not} produced by Data Exchange, whose timetable export is the class-assignment
interchange format and cannot be loaded by the solver. Two settings matter when producing it:
\emph{Student Course Demands} must be \emph{Enrolled Student Course Demands}, or the student block
is empty, and \texttt{unitime.solution.export.names} must be \texttt{true}, or class and room names
are stripped.

\subsection{Derived artefacts}

\begin{table}[H]
\centering
\caption{Artefacts derived from the Spr6 session.}
\label{tab:spr6art}
\footnotesize
\begin{tabular}{@{}L{5.4cm}L{2.0cm}L{7.0cm}@{}}
\toprule
\textbf{Artefact} & \textbf{Type} & \textbf{Role} \\
\midrule
\texttt{rl/instances/spr2026.xml}            & solver XML & the instance above; input to every sweep run \\
\texttt{rl/instances/spr2026-nostudents.xml} & solver XML & the same export before student injection \\
\texttt{rl/instances/spr2026v6\_mihir.xml},
  \texttt{spr2026v8\_mihir.xml}              & solver XML & later exports of the same session \\
\texttt{solutions/spr2026.kpi.json}          & KPI report & scores the instance: 526/990, hard feasibility 92.15 \\
\texttt{solutions/COEPSpr62026\_v2.csv}      & CSV        & exported timetable, scored in \texttt{kpi\_report.json} \\
\texttt{solutions/course-solution\_weight5.csv},
  \texttt{\_weight20.csv}                    & CSV        & the UniTime weight experiments of Section~\ref{sec:experiments} \\
\texttt{solutions/spr2026\_10m\_w5.csv}      & CSV        & ten-minute run at weight 5 \\
\texttt{rl/runs/sweep/}                      & run tree   & the controlled weight sweep: 7 configurations $\times$ 3 seeds \\
\texttt{rl/runs/bakeoff-spr2026\_10m\_w5/}   & run tree   & arms B0, B1, B2 $\times$ 10 seeds on an already-solved instance \\
\texttt{rl/runs/sweep/w5/.../solution.xml}   & solver XML & scored independently in \texttt{solutions/solution.kpi.json} \\
\bottomrule
\end{tabular}
\end{table}

\subsection{Why 990 and not 1{,}290}

The dataset of this period declared 1{,}290 simulated students. The solver export carries 990. The
difference is not attrition: 990 is the number of students for whom the solver had loaded
enrolments at the moment of export. A conflict rate computed over 990 and one computed over
1{,}290 are therefore measurements of different populations, which is why
Table~\ref{tab:crossrun} records the denominator for every run.
